% =============================================================================
% CAPÍTULO 1 — INTRODUCCIÓN
% =============================================================================
\chapter{Introducción}
\label{ch:introduccion}
\resumenCapitulo{Este capítulo presenta el propósito, el alcance y el público objetivo
del manual, y describe las convenciones tipográficas y visuales que se emplean a lo
largo del documento para guiar a la persona usuaria paso a paso.}

% -----------------------------------------------------------------------------
\section{Propósito del Documento}
\label{sec:proposito}

El presente \textbf{Manual del Sistema para Perfiles} tiene como propósito guiar a las
personas usuarias de la plataforma \textbf{EthosHub} en el uso correcto y completo de
todas las funciones disponibles según su perfil. EthosHub es un sistema de gestión de
talento y reclutamiento basado en valores, en el que las personas profesionales
construyen su portafolio digital, las personas reclutadoras buscan y evalúan talento, y
el personal administrador supervisa la operación de la plataforma.

A diferencia del manual técnico —orientado a la arquitectura y al mantenimiento del
software— este documento describe el \textit{qué} y el \textit{cómo} desde la
perspectiva de quien opera el sistema: dónde hacer clic, qué datos introducir y qué
resultado esperar después de cada acción.

\begin{AvisoNota}
Cada procedimiento de este manual se acompaña de capturas de pantalla reales del
sistema. Las instrucciones están redactadas en orden secuencial: basta con seguir los
pasos en el orden presentado para completar cada tarea con éxito.
\end{AvisoNota}

\paragraph{Conformidad normativa.} Este manual se elabora conforme al estándar
\textbf{ISO/IEC/IEEE 26514:2022}, \textit{Systems and software engineering --- Design
and development of information for users}, que establece los criterios de estructura,
contenido y presentación de la documentación de usuario.

% -----------------------------------------------------------------------------
\section{Alcance del Sistema}
\label{sec:alcance}

El manual cubre la totalidad de las funciones operativas de la plataforma EthosHub,
organizadas por perfil de usuario:

\begin{itemize}[noitemsep]
    \item \textbf{Perfil Profesional} (capítulo~\ref{ch:profesional}): construcción y
    gestión del portafolio (proyectos, experiencia laboral, formación académica y
    habilidades), generación de currículum con CV Studio, y configuración del perfil
    público.
    \item \textbf{Perfil Reclutador} (capítulo~\ref{ch:reclutador}): panel de
    reclutamiento, búsqueda inteligente y filtrada de talento, rankings, gestión del
    \textit{pipeline} de selección y comunicación directa con los candidatos.
    \item \textbf{Perfil Administrador} (capítulo~\ref{ch:administrador}): panel de
    administración (\textit{backoffice}), gestión de cuentas, moderación de contenido,
    administración del repositorio de habilidades, métricas globales y comunicaciones.
\end{itemize}

Asimismo, el manual incluye la gestión de acceso (registro, inicio de sesión y
recuperación de credenciales), el uso de la plataforma desde dispositivos móviles
(diseño adaptable), la resolución de problemas frecuentes y un glosario de términos.

\paragraph{Fuera de alcance.} Quedan excluidos de este documento los procedimientos de
instalación y despliegue (cubiertos por el \textit{Manual de Instalación}) y el diseño
arquitectónico interno del software (cubierto por el \textit{Manual Técnico}).

% -----------------------------------------------------------------------------
\section{Público Objetivo}
\label{sec:publico}

Este manual está dirigido a tres audiencias, correspondientes a los tres perfiles del
sistema. No se requieren conocimientos técnicos de programación; únicamente se asume
familiaridad básica con el uso de un navegador web y de formularios en línea.

\begin{TablaInfo}{Perfil}{Descripción de la persona usuaria}
\textbf{Profesional} & Persona que desea publicar su portafolio, demostrar sus
habilidades y proyectos, y hacerse visible ante reclutadores. \\ \hline
\rowcolor{grisTecnico}
\textbf{Reclutador} & Persona que busca, evalúa y contacta talento para procesos de
selección, gestionando candidatos a lo largo de un \textit{pipeline}. \\ \hline
\textbf{Administrador} & Persona responsable de la operación, la moderación y la
supervisión global de la plataforma. \\ \hline
\end{TablaInfo}

% -----------------------------------------------------------------------------
\section{Convenciones de la Guía y Formato Visual}
\label{sec:convenciones}

Para facilitar la lectura y la ejecución de los procedimientos, este manual emplea un
conjunto uniforme de convenciones tipográficas y visuales.

\subsection{Resaltados de Texto y Nomenclatura de Botones}
\label{subsec:nomenclatura}

\begin{itemize}[noitemsep]
    \item Los \textbf{botones} y controles sobre los que se debe hacer clic se
    representan así: \boton{Guardar}, \boton{Crear proyecto}, \boton{Cancelar}.
    \item Los \textbf{campos de formulario} y etiquetas de la interfaz se muestran en
    cursiva: \campo{Nombre}, \campo{Correo electrónico}, \campo{Nivel profesional}.
    \item Las \textbf{rutas de navegación} dentro de la aplicación se indican con el
    símbolo «\textgreater»: \rutanav{Configuración \textgreater\ Seguridad \textgreater\ Cambiar contraseña}.
    \item Los \textbf{valores literales}, direcciones y comandos se muestran en
    tipografía monoespaciada, por ejemplo \comando{https://ethoshub.com}.
    \item Los \textbf{resultados esperados} de una acción se destacan en verde:
    \resultado{Proyecto creado correctamente}.
\end{itemize}

Además, el manual utiliza tres tipos de avisos destacados:

\begin{AvisoNota}
Aporta consejos de productividad e información útil para aprovechar mejor una función.
\end{AvisoNota}

\begin{AvisoPrecaucion}
Advierte sobre acciones que pueden producir un resultado no deseado si no se realizan
con cuidado.
\end{AvisoPrecaucion}

\begin{AvisoAdvertencia}
Señala acciones irreversibles o que pueden ocasionar la pérdida de datos, como la
eliminación de una cuenta.
\end{AvisoAdvertencia}

\subsection{Interpretación de Capturas de Pantalla y Simbología}
\label{subsec:capturas}

Cada procedimiento se ilustra con capturas de pantalla numeradas como
\textbf{Figura~\textit{capítulo.número}}. Las capturas se presentan inmediatamente
después del paso que documentan, y su leyenda describe con precisión el control o la
acción resaltada en la imagen.

\begin{itemize}[noitemsep]
    \item Las imágenes se muestran a un tamaño uniforme para conservar la legibilidad y
    se disponen como máximo dos por página.
    \item Cuando un procedimiento depende de un paso anterior, la leyenda lo indica
    explícitamente (por ejemplo, «tras pulsar \boton{Continuar}\ldots»).
    \item Los iconos del sistema se describen por su forma (icono de lápiz para editar,
    icono de papelera para eliminar, icono de estrella para destacar, etc.).
\end{itemize}

\begin{NotaTecnica}
Los datos que aparecen en las capturas (nombres, correos, proyectos) son ejemplos
ilustrativos. La información real dependerá de la cuenta y del contenido de cada
persona usuaria.
\end{NotaTecnica}
